% Generated draft: requires mathematical and source review.
\documentclass[11pt]{article}
\usepackage[margin=25mm]{geometry}
\usepackage{fontspec}
\setmainfont{Latin Modern Roman}
\usepackage{amsmath,amssymb,mathtools,unicode-math}
\setmathfont{Latin Modern Math}
\usepackage{enumitem,needspace,xurl,hyperref}
\setlength{\parindent}{0pt}
\setlength{\parskip}{5pt}
\setlength{\emergencystretch}{3em}
\newcommand{\N}{\mathbb{N}}
\newcommand{\Z}{\mathbb{Z}}
\newcommand{\Q}{\mathbb{Q}}
\newcommand{\R}{\mathbb{R}}
\newcommand{\C}{\mathbb{C}}
\newcommand{\F}{\mathbb{F}}
\newcommand{\A}{\mathbb{A}}
\newcommand{\PP}{\mathbb{P}}
\newcommand{\E}{\mathbb{E}}
\newcommand{\eps}{\varepsilon}
\newcommand{\dd}{\,\mathrm{d}}
\begin{document}
\section*{UnsolvedMath ID: 800481}
\textbf{Source problem code:} catalogue entry
\subsection{Definitions and mathematical statement}
Let \(P\) be a polygon, let \(S(P)\) be its finite set of sides, and let \(\lambda:S(P)\to\mathcal A\) assign symbols from a finite alphabet \(\mathcal A\) to its sides. Let \(\mathcal T(P)\) denote the intended collection of polygonal symplectic billiard trajectories admitting the bi-infinite side coding described in the source. For \(\gamma\in\mathcal T(P)\), write its code as
\[
 c(\gamma)=(c_i(\gamma))_{i\in\Z}\in\mathcal A^{\Z},
\]
where \(c_i(\gamma)\) is the label of the side recorded at index \(i\). For every integer \(n\geq1\), define
\[
 \mathcal L_{P,\lambda}(n)
 =\bigl\{(c_k(\gamma),\ldots,c_{k+n-1}(\gamma)):
       \gamma\in\mathcal T(P),\ k\in\Z\bigr\},
 \qquad
 p_{P,\lambda}(n)=\#\mathcal L_{P,\lambda}(n).
\]
Thus \(p_{P,\lambda}(n)\) counts distinct contiguous words of length \(n\) occurring across the entire trajectory collection, rather than counting their occurrences.

\textbf{Problem (Serge Tabachnikov, Problem 3).} Study the complexity function \(p_{P,\lambda}(n)\) of polygonal symplectic billiards, including its growth as \(n\to\infty\). The supplied source poses a research problem and specifies no particular conjectured bound or classification.

The source does not define the polygonal symplectic billiard dynamics, specify the permitted class of polygons, or explain how singular trajectories and trajectories without bi-infinite continuation are treated. Consequently, \(\mathcal T(P)\) above is a placeholder for the source's intended trajectory collection, not a completed definition of that dynamics. The source also does not explicitly require distinct symbols for distinct sides; the notation retains the labeling \(\lambda\) because the complexity can depend on it. These qualifications must be resolved before asserting a statement for a specified class of polygons and trajectories.

For context, the supplied record states that conventional polygonal billiards have subexponential complexity and that rational polygonal billiards have polynomial, in fact cubic, growth; here rational means that all interior angles are rational multiples of \(\pi\). It also states that polygonal outer billiards have polynomial complexity growth. These are supplied background claims, with their hypotheses and bibliographic support not fully specified in the record, and are not conclusions of this draft. Finally, the source reports examples of polygonal symplectic billiards for which all trajectories are periodic and describes their complexity as having a zero growth rate. It does not define that growth rate or provide a uniform period bound, so this observation is retained as a qualified source report without identifying it with a particular asymptotic formula.
\subsection{Short English statement}
Serge Tabachnikov asks how the number of distinct length-\(n\) side-label words in polygonal symplectic billiard trajectories grows with \(n\). The symbolic counting function is defined above, but the supplied record leaves the precise dynamics, admissible polygons, singularity conventions, and meaning of zero growth rate unresolved.
\subsection{Sources}
\begin{itemize}
\item Local UnsolvedMath registry, permanent ID \(800481\), exact source problem code \texttt{catalogue entry}; supplied title: \emph{Billiards and Quantitative Symplectic Geometry --- Serge Tabachnikov --- Problem 3}. The registry labels the problem open; this is supplied status metadata, not an independently verified assessment.
\item Albach, Barraud et al., \emph{Open Problems in Billiards and Quantitative Symplectic Geometry} (2026), source item Serge Tabachnikov, Problem 3. Supplied source URL: \url{https://arxiv.org/abs/2602.12896}. Contributor: Serge Tabachnikov. The background records access on 2026-10-16 and extraction by \texttt{arxiv\_latex}; these are supplied provenance statements, not a claim of access or verification for this draft.
\item Reference label [AT] is supplied for the introduction of symplectic billiards; labels [ABSST] and [ALW] are supplied for studies of polygonal symplectic billiards. The record provides no full bibliographic entries or individual links for these labels.
\item Reference label [Ka] is supplied as an early reference for the complexity literature on conventional polygonal billiards. Reference label [GT] is supplied for polynomial complexity growth in polygonal outer billiards. The record provides no full bibliographic entries or individual links for these labels.
\end{itemize}
\end{document}
